\section{Ước lượng tham số: vòng lặp ALS}
\label{sec:estimation}

\subsection{Bài toán tối ưu bằng trực giác}

Ta có dữ liệu $\mathcal{X}^{c}$ và một mô hình gồm nhiều khối: lõi trend, hệ reaction--diffusion,
hệ số từ khóa/địa điểm, hệ số mùa vụ và outlier. Nếu tối ưu tất cả cùng lúc, bài toán vừa phi
tuyến vừa nhiều tham số. Paper dùng \textbf{ALS} (alternating least squares): giữ phần lớn tham
số cố định, cập nhật một khối; lặp lại đến khi hội tụ.

\begin{keypoint}
ALS là chiến lược ``mỗi lần chỉ sửa một phần, nhưng sửa đến khi phần còn lại cho phép''. Nó không
đảm bảo tìm cực trị toàn cục cho bài toán không lồi; paper dùng nó như một thủ tục thực dụng và
báo cáo kết quả thực nghiệm.
\end{keypoint}

Để dễ theo dõi, tài liệu gọi sai số bình phương của một mô hình là
\begin{equation}
  \mathcal{L}(\Theta)
  =\left\|\mathcal{X}^{c}
    -\widehat{\mathcal{X}}_d-\widehat{\mathcal{X}}_s-\widehat{\mathcal{X}}_o
   \right\|_F^2.
  \label{eq:als-loss}
\end{equation}
Paper không đặt tên ký hiệu này trong mục 4.1, nhưng mọi bước cập nhật dưới đây đều nhằm giảm
sai số đó (hoặc giảm chi phí MDL ở bước outlier). Chuẩn Frobenius $\|\cdot\|_F$ chỉ là căn bậc hai
tổng bình phương mọi phần tử, tức một cách đo ``mô hình lệch dữ liệu bao nhiêu''.

\subsection{Khởi tạo: STL, PARAFAC và NTD}

Thuật toán ModelEstimation cần điểm bắt đầu. Paper mô tả ba bước khởi tạo:
\begin{enumerate}
  \item Dùng \textbf{STL} (seasonal-trend decomposition using LOESS) để tách cửa sổ hiện tại
        thành trend, seasonal và residual. STL là phương pháp thống kê cổ điển, mạnh với chuỗi
        có mùa vụ; nó cho ta xấp xỉ ban đầu của $\widehat{\mathcal{X}}_d$ và
        $\widehat{\mathcal{X}}_s$.
  \item Phân rã $\widehat{\mathcal{X}}_s$ bằng \textbf{PARAFAC} để có các hệ số mùa vụ ban đầu
        $\mathbf{S}^{(\mathrm{time})},\mathbf{S}^{(\mathrm{key})},\mathbf{S}^{(\mathrm{loc})}$.
  \item Áp dụng \textbf{Nonnegative Tucker Decomposition (NTD)} cho $\widehat{\mathcal{X}}_d$ để
        lấy $\mathbf{W}^{(\mathrm{key})},\mathbf{W}^{(\mathrm{loc})}$ ban đầu.
\end{enumerate}
Sau đó $\widehat{\mathcal{X}}_o$ được đặt bằng 0: giả định ban đầu là chưa có outlier. Việc khởi
tạo không phải chi tiết phụ; nó quyết định vòng ALS có rơi vào một nghiệm xấu hay không, và paper
ghi rõ cách làm trong Appendix A.1.

\subsection{Cập nhật hệ reaction--diffusion}

Giữ $\mathbf{W}^{(\mathrm{key})},\mathbf{W}^{(\mathrm{loc})}$ và các thành phần khác cố định,
paper giải bài toán bình phương tối thiểu phi tuyến
\begin{equation}
  \{\theta,w_0\}
  =\arg\min_{\theta',w_0'}
   \left\|
     \mathcal{X}^{c}-\widehat{\mathcal{X}}_s-\widehat{\mathcal{X}}_o
     -f_\theta(\theta',w_0',L_c)
       \times_{\mathrm{key}}\mathbf{W}^{(\mathrm{key})}
       \times_{\mathrm{loc}}\mathbf{W}^{(\mathrm{loc})}
   \right\|_F^2,
  \label{eq:lm-estimation}
\end{equation}
đây là Eq.~(5) trong paper. Ký hiệu $\arg\min$ nghĩa là ``chọn bộ tham số làm biểu thức nhỏ nhất''.
Vì $f_\theta$ chứa nghiệm của ODE, đạo hàm theo $\theta$ không có dạng đóng đơn giản. Paper dùng
\textbf{Levenberg--Marquardt (LM)}, một phương pháp số cho bình phương tối thiểu phi tuyến, để cập
nhật đồng thời các tham số $\mathbf{A},\mathcal{D}$ và trạng thái ban đầu $w_0$.

Ta không cần tự viết LM để hiểu paper, nhưng cần nhớ ba điều: (i) đây là bước đắt nhất về mặt
tính toán; (ii) nó cần một bộ giải ODE; (iii) kết quả phụ thuộc vào điểm khởi tạo và có thể chỉ là
cực trị địa phương.

\subsection{Cập nhật hệ số trend bằng multiplicative update}

Khi đã có lõi $\mathcal{W}^{(\mathrm{core})}$, paper cập nhật từng hệ số trend. Với
``mode'' là key hoặc loc, đặt
\begin{align}
  \mathbf{P}^{(\mathrm{mode})}
  &=\max\!\left(\epsilon,\;
    \mathbf{X}_d^{(\mathrm{mode})}\cdot\mathbf{G}^{(\mathrm{mode})}
   \right), \label{eq:w-p}\\
  \mathbf{Q}^{(\mathrm{mode})}
  &=\max\!\left(\epsilon,\;
    \mathbf{W}^{(\mathrm{mode})}\cdot\mathbf{G}^{(\mathrm{mode})}
    \cdot(\mathbf{G}^{(\mathrm{mode})})^{\mathsf T}
   \right), \label{eq:w-q}\\
  \mathbf{W}^{(\mathrm{mode})}
  &\leftarrow
    \mathbf{W}^{(\mathrm{mode})}
    \otimes\mathbf{P}^{(\mathrm{mode})}
    \oslash\mathbf{Q}^{(\mathrm{mode})}.
  \label{eq:multiplicative}
\end{align}
Đây là Eq.~(6) trong paper. Các ký hiệu:
\begin{itemize}
  \item $\mathbf{X}_d^{(\mathrm{mode})}$ là tensor residual
        $\mathcal{X}^{c}-\widehat{\mathcal{X}}_s-\widehat{\mathcal{X}}_o$ sau khi \emph{unfold}
        theo mode tương ứng (Eq.~7);
  \item $\mathbf{G}^{(\mathrm{mode})}$ là phần còn lại của phép chiếu, cũng được unfold;
  \item $\otimes$ là nhân từng phần tử (Hadamard), $\oslash$ là chia từng phần tử;
  \item $\epsilon$ là hằng số dương rất nhỏ để tránh chia cho 0.
\end{itemize}
Vì $\mathbf{W}$ ban đầu không âm và phép cập nhật chỉ nhân/chia các đại lượng dương, hệ số vẫn
không âm. Đây là lý do ``multiplicative'' quan trọng: nó vừa tối ưu vừa giữ ràng buộc.

\begin{plainwords}
Công thức~(6) nghe phức tạp, nhưng ý tưởng là: so sánh phần dữ liệu còn thiếu với phần mô hình
dự đoán, rồi tăng/giảm từng trọng số theo tỉ lệ. Không có phép trừ tự do nào có thể đẩy trọng số
xuống âm.
\end{plainwords}

\subsection{Cập nhật hệ số mùa vụ}

Với phần mùa vụ, paper giữ các thành phần khác cố định và dùng một bước ALS có dạng đóng. Với mỗi
mode trong $\{\mathrm{time},\mathrm{key},\mathrm{loc}\}$, đặt
\begin{align}
  \mathbf{U}^{(\mathrm{mode})}
  &=\odot_{i\neq\mathrm{mode}}\mathbf{S}^{(i)},&
  \mathbf{V}^{(\mathrm{mode})}
  &=\otimes_{i\neq\mathrm{mode}}
    (\mathbf{S}^{(i)})^{\mathsf T}\mathbf{S}^{(i)},\\
  \mathbf{S}^{(\mathrm{mode})}
  &\leftarrow
    \mathbf{X}_s^{(\mathrm{mode})}
    \mathbf{U}^{(\mathrm{mode})}
    (\mathbf{V}^{(\mathrm{mode})})^{\dagger}.
  \label{eq:seasonal-update}
\end{align}
Đây là Eq.~(8); $\odot$ và $\otimes$ lần lượt là tích Khatri--Rao và tích từng phần tử,
$\dagger$ là giả nghịch đảo Moore--Penrose, còn $\mathbf{X}_s^{(\mathrm{mode})}$ là residual
$\mathcal{X}^{c}-\widehat{\mathcal{X}}_d-\widehat{\mathcal{X}}_o$ được unfold (Eq.~9). Về bản
chất, đây là nghiệm bình phương tối thiểu tuyến tính khi giữ hai mode còn lại; không cần chứng
minh lại trong tài liệu này.

\subsection{Cập nhật outlier bằng sparsification}

Sau khi cập nhật trend và mùa vụ, residual còn lại là
$\mathcal{R}=\mathcal{X}^{c}-\widehat{\mathcal{X}}_d-\widehat{\mathcal{X}}_s$. Paper không dùng
một ngưỡng thủ công kiểu ``lớn hơn 3 thì là outlier''. Nó xét từng phần tử của $\mathcal{R}$ và
quyết định giữ hay đặt về 0 sao cho tổng chi phí MDL nhỏ nhất. Chi tiết chi phí ở
Chương~\ref{sec:mdl-streaming}; ở đây chỉ cần hiểu kết quả là một tensor thưa.

\subsection{Toàn bộ vòng lặp và độ phức tạp}

\begin{tcolorbox}[enhanced,breakable,colback=white,colframe=DTNavy,arc=1.5mm,
  title={\sffamily\bfseries Algorithm 2 --- ModelEstimation (tóm tắt)},fonttitle=\small,
  coltitle=white,colbacktitle=DTNavy,left=3mm,right=3mm,top=2mm,bottom=2mm]
\textbf{Input:} $\mathcal{X}^{c},d_k,d_l,d_s$.\quad
\textbf{Output:} $\Theta=\{\theta_d,\theta_s,\widehat{\mathcal{X}}_o\}$.
\begin{enumerate}[leftmargin=1.5em]
  \item Khởi tạo $\theta_d,\theta_s$ bằng STL + PARAFAC + NTD; đặt $\widehat{\mathcal{X}}_o=0$.
  \item \textbf{repeat}
  \item \quad Cập nhật $\theta=\{\mathbf{A},\mathcal{D}\}$ và $w_0$ bằng LM (Eq.~5).
  \item \quad Sinh $\mathcal{W}^{(\mathrm{core})}=f_\theta(\theta,w_0,L_c)$.
  \item \quad Cập nhật $\mathbf{W}^{(\mathrm{key})},\mathbf{W}^{(\mathrm{loc})}$ (Eq.~6).
  \item \quad Cập nhật $\mathbf{S}^{(\mathrm{time})},\mathbf{S}^{(\mathrm{key})},
        \mathbf{S}^{(\mathrm{loc})}$ (Eq.~8).
  \item \quad Sparsify residual để cập nhật $\widehat{\mathcal{X}}_o$.
  \item \textbf{until} hội tụ.
\end{enumerate}
\end{tcolorbox}

Lemma~1 (trong §4.1; HTML đặt nhãn là Lemma~4.1) phát biểu độ phức tạp thời gian của ModelEstimation là
\[
  O(d_k d_l^2+kd_k+ld_l).
\]
Con số này đến từ $d_k d_l^2$ tham số của hệ reaction--diffusion và chi phí cập nhật hai bảng hệ
số. Tuy nhiên, biểu thức đã bỏ qua số vòng lặp ALS, chi phí giải ODE, chi phí cập nhật mùa vụ và
outlier. Vì vậy, hãy đọc đây là độ phức tạp của \emph{một vòng cập nhật chính} theo cách paper
trình bày, không phải bảo đảm cho toàn bộ pipeline.

\begin{commonmistake}
``Hội tụ'' không có nghĩa tìm được nghiệm tốt nhất toàn cục. ALS trên bài toán không lồi có thể
dừng ở cực trị địa phương; paper không chứng minh tính tối ưu toàn cục. Ngoài ra, $\epsilon$ không
phải tham số mô hình cần học; nó chỉ bảo vệ phép chia.
\end{commonmistake}

\begin{selfcheck}
\begin{enumerate}
  \item Vì sao bước cập nhật $\theta$ cần LM thay vì một nghiệm đóng?
  \item Phép cập nhật nhân $\mathbf{W}\otimes\mathbf{P}\oslash\mathbf{Q}$ giữ được tính không âm nhờ đâu?
  \item Kể tên ba bước khởi tạo và phương pháp tương ứng.
  \item Độ phức tạp $O(d_kd_l^2+kd_k+ld_l)$ bỏ qua những chi phí nào?
\end{enumerate}
\end{selfcheck}
